\documentclass[10pt,a4paper]{amsart}
\usepackage[margin=19mm]{geometry}
\usepackage{amsmath,amssymb,mathtools}
\usepackage[scr=rsfs,cal=euler]{mathalfa}
\usepackage{xfrac}
\usepackage[unicode,hidelinks,bookmarks=false]{hyperref}
\newcommand{\ie}{i.e.\ }
\newcommand{\cf}{cf.\ }
\newcommand{\resp}{respectively\ }
\newcommand{\Subsection}{\subsection}
\providecommand{\nobreakdash}{\nobreak\hbox{-}}
\setlength{\emergencystretch}{2em}
\multlinegap0pt
\usepackage{bm}
\usepackage{bbm}
\usepackage{dsfont}
\usepackage{xspace}
\usepackage{cancel}
\usepackage{mathtools}
\usepackage{amsthm}
\makeatletter
\@ifundefined{proposition}{\newtheorem{proposition}{Proposition}}{}
\@ifundefined{theorem}{\newtheorem{theorem}{Theorem}}{}
\makeatother
\newcommand{\case}[1]{\paragraph*{Case #1:}}
\newcommand{\irr}{\operatorname{irr}}
\title[Bounds on Trees with Topological Indices Among Degree Sequen]{Bounds on Trees with Topological Indices Among Degree Sequence}
\author{Jasem Hamoud, Alexei Belov-Kanel, Duaa Abdullah}
\date{Source: arXiv:2505.12518v2 (CC BY 4.0)}
\begin{document}
\maketitle

\noindent\emph{Source and licence.} Bounds on Trees with Topological Indices Among Degree Sequence. Version arXiv:2505.12518v2 (CC BY 4.0), distributed under the Creative Commons Attribution 4.0 International licence, as recorded in the arXiv licence metadata for this exact version and shown on the abstract page https://arxiv.org/abs/2505.12518v2. Full rights notice: licenses/arxiv-2505.12518-CC-BY-4.0.txt.\par\smallskip

\noindent\textit{Editorial scope.} Counted unit: the Theorem whose source label is \texttt{thmsign1} in the article (source0.tex lines 579--584, source label \texttt{thmsign1}), delivered together with the article's own complete original proof of it (source0.tex lines 585--647, one \texttt{proof} environment of 63 lines). No mathematics is written, repaired or re-derived: statement and proof are the article's own text line for line. Required context retained uncounted: prozagn2 (lines 261--264); mainalb4 (lines 372--377); thmalbn1 (lines 426--431). Every definition, display and label of those lines is retained unchanged. Omitted from this package and not redistributed: 1--260, 265--371, 378--425, 432--578, 648--810. Nothing inside a retained excerpt is omitted or altered: every retained body is delivered exactly as the article's lines are written, so no omission receipt of any kind accompanies this package. Citations: the retained excerpts carry no citation command at all, so no bibliography entry of the article is transcribed and no reference is redistributed. Reference adaptation: none was needed and none was made; every reference inside the delivered unit resolves inside it, so no unresolved reference or citation is produced. Printed slips kept literal: no printed word, symbol, hypothesis or number of the counted statement or of its proof was altered. Layout and class: the article's own class is \texttt{amsart} with packages including graphicx, amssymb, stmaryrd, enumitem, lmodern, geometry, tikz, float, url, mathtools, amsmath, amsfonts, amsthm, mathrsfs; the wrapper is the lane's pinned \texttt{amsart} editorial wrapper at 10pt on A4, which restates the theorem-like environments the retained text needs and adds the article's own macro definitions that the retained text needs (\texttt{\textbackslash case}, \texttt{\textbackslash irr}), with extra wrapper packages (bm, bbm, dsfont, xspace, cancel, mathtools). The delivered numbering is the unit's own. Assets: no retained span contains a figure environment, a graphics-inclusion command, a PDF-page inclusion, a table or a TikZ picture; no image file is copied into this package, the package declares no asset dependency and no image file is redistributed with it. Licence: version arXiv:2505.12518v2 of this article is distributed under the Creative Commons Attribution 4.0 International licence, as recorded in the arXiv licence metadata for this exact version and shown on the abstract page https://arxiv.org/abs/2505.12518v2. A full rights notice is delivered at licenses/arxiv-2505.12518-CC-BY-4.0.txt. Characters: every retained excerpt is pure ASCII, so the delivered engine, XeLaTeX with its default Latin Modern text font, reports no missing character.
\begin{proposition}~\label{prozagn2}
Let $T=(V,E)$ be a tree where vertices set $V=\{v_1,\dots,v_i\}$ with $n=|V|$ vertices and edges set $E=\{e_1,\dots,e_j\}$ with $m=|E|$ edges, let $\mathscr{D}=(d_1,\dots,d_n)$ be a degree sequence, where $d_n\geqslant d_{n-1}\geqslant \dots \geqslant d_1$, let $M_1(T)$ be the first Zagreb index with maximum degree $\Delta$, then Albertson index satisfying 
\[\irr(T)\geqslant 2\sum_{v_i\in V(T)}\binom{d_i}{2}+\frac{4m}{n-1}-\Delta(\Delta-1).\]
\end{proposition}

\begin{theorem}~\label{mainalb4}
Let $T=(V,E)$ be a regular tree with maximum degree $\Delta$ and minmum  degree $\delta$, then  an inequality involving the irregularity of $T$ by Albertson index satisfying: 
    \[
    \irr(T) \leq \delta \left( \frac{2\delta}{\Delta - 1} + \Delta - 1 \right) + \Delta (\Delta - 1).
    \]
\end{theorem}

\begin{theorem}~\label{thmalbn1}
Let $T$ be a tree with $n$ vertices and $m$ edges with maximum degree $\Delta$ and minmum  degree $\delta$, let $M_1(T)$ be the first Zagreb index, then inequality of Albertson index $\irr$ satisfying 
\[
\irr(T) \geq M_1(T)  - \frac{2mn^2(\Delta - 1)+\delta-1}{n+\Delta}.
\]
\end{theorem}

\begin{theorem}~\label{thmsign1}
Let $T$ be a tree with $n$ vertices and $m$ edges with maximum degree $\Delta$ and minmum  degree $\delta$, let $M_1(T)$ be the first Zagreb index, then inequality of Sigma index $\irr$ satisfying 
\[
\sigma(T) \geq M_1(T)  - \frac{2mn^2(\Delta - 1)+\delta-1}{n+\Delta}.
\]
\end{theorem}

\begin{proof}
Let $T$ be a regular tree, $S_n$ a star tree of order $n\geq 3$, $P_n$ path of length $n$ at least degree 1. In fact, from Theorem~\ref{thmalbn1}, introduced the bound 
\begin{equation}~\label{tessign1}
\irr(T) \geq M_1(T)  - \frac{2mn^2(\Delta - 1)+\delta-1}{n+\Delta}.
\end{equation}
This inequality provides a lower bound on the Albertson index of a tree $T$ in terms of the first Zagreb index $M_1(T)$, where $M_1(T)$ relate different structural measures and to estimate irregularity. 
\case{a} Let $S_n$ be a star tree of order $n\geq 3$. Since $\sigma(T)=\sum_{uv\in E}(d_u-d_v)^2$, and $\irr(T)=\sum_{uv\in E}\lvert d_u-d_v\rvert$, then when $T$ is a star tree, we have the bound satisfying
\begin{equation}~\label{sigmaboundn1}
    \sigma(S_n) \geq \Delta^3 + \Delta^2 - \frac{2 \Delta (\Delta + 1)^2 (\Delta - 1)}{2\Delta^2 + 1}.
\end{equation}
\case{1} If $\Delta=n-1$, Theorem~\ref{mainalb4} established that for a star tree $\irr(S_n)=(n-1)(n-2)$, then we find $M_1(S_n)=n(n-1)$, and when $n\geq 3$ we have $\sigma(S_n)=(n-1)(n-2)$, then the lower bound is
\begin{equation}~\label{eqqsigmasatn1}
\sigma(S_n) \geq (n-1)^3+2n^2.
\end{equation}
Hence, we have
\begin{align*}
 \sigma(S_n)-M_1(S_n) & \geq (n-1)^3+2n^2-n(n-1)\\
 &=n^3 - 2 n^2 + 4 n - 1>0.
\end{align*}
Thus, bound~(\ref{sigmaboundn1}) observe that when $\Delta(\Delta-1)\geq 0$, we have $\sigma(S_n) \geq 2(n-1)n^2(\Delta - 1)$ is the term of inequality~(\ref{sigmaboundn2}), then $\sigma(S_n) \geq 2(n-1)n^2(\Delta - 1)+\delta-1$, when $n+\Delta>0$ the inequality~(\ref{sigmaboundn2} valid. When $\sigma(S_n)-M_1(S_n)>0$ clearly. 
\begin{equation}~\label{sigmaboundn2}
     \sigma(S_n) \geq M_1(S_n)  - \frac{2mn^2(\Delta - 1)+\delta-1}{n+\Delta}.
\end{equation}
\case{2} In this case, if $\Delta=n+1$, then from~(\ref{eqqsigmasatn1}) the upper bound is 
\begin{equation}~\label{eqqsigmasatn2}
\sigma(S_n) \geq (n+1)^3+2n^2.
\end{equation}
Thus, $\sigma(S_n)-M_1(S_n)>0$, then~(\ref{sigmaboundn2}) is valid.
\case{b} If $T$ be a path denote by $P_n$, then according to~(\ref{tessign1}) we find that $ M_1(P_n) = 4n - 6$, according to Proposition~\ref{prozagn2} we have $\sigma(T)\geq M_1(T)+n\Delta^2(\Delta-1)+m$, then we provide an upper bound $\sigma(T)\geq M_1(T)+2mn^2\Delta^2(\Delta-1)$, so that we know $m=n-1$ for any tree, then $\sigma(T)\geq M_1(T)+2(n-1)n^2\Delta^2(\Delta-1)$ holds for $k>0$ the bound explicitly
\begin{equation}~\label{eqqsigmasatn3}
\sigma(T)\geq \Delta(\Delta-1)+k.\frac{2(n-1)n^2\Delta^2(\Delta-1)}{n+\Delta}.
\end{equation}
\case{1} If $k=\delta$, then from~(\ref{eqqsigmasatn3}) involves the Sigma index $\sigma(T)$ related to the path $P_n$, where Sigma index measures certain structural properties related to vertex degrees and adjacency, then the bound is 
\begin{equation}~\label{eqqsigmasatn4}
\sigma(P_n)\geq \Delta(\Delta-1)\left (1+\delta\frac{2(n-1)n^2\Delta}{n+\Delta}\right).
\end{equation}
Actually, from~(\ref{eqqsigmasatn4}), for $\Delta=n-1$, we define the lower bound
\begin{equation}~\label{sigmalowern1}
    (n - 1)(n - 2) \left( 1 + \delta \cdot \frac{2 n^{2} (n - 1)^{2}}{2 n - 1} \right).
\end{equation}
Then, for $\Delta=n+1$, we find that the upper bound 
\begin{equation}~\label{sigmauppern1}
    n \left( n + 1 + \delta \cdot \frac{2 (n - 1) n^{2} (n + 1)^{2}}{2 n + 1} \right).
\end{equation}
Both bounds~(\ref{sigmalowern1}),(\ref{sigmauppern1}) increase with the minimum degree $\delta$, reflecting that maximum degree $\Delta=n+1$ and minimum degree to increase the value of the respective tree $T$. This bounds~(\ref{sigmalowern1}),(\ref{sigmauppern1}) declare $\sigma(P_n)-M_1(P_n)>0$. 
\case{2} If $k=\Delta$, then from~(\ref{eqqsigmasatn3}) holds to 
\begin{equation}~\label{eqqsigmasatn5}
    \Delta (\Delta - 1) \left( 1 + \frac{2 (n - 1) n^{2} \Delta^{2}}{n + \Delta} \right).
\end{equation}
Hence, the lower bound when $\Delta=n-1$, both of term $2n-1$, $2n+1$ are close for large $n$, then $\sigma(T)= \Delta (\Delta - 1)$ satisfying by considering to~(\ref{eqqsigmasatn1}) as 
\begin{equation}~\label{eqqsigmalowern2}
   (n-1)(n-2)\left( 1 + \frac{2 n^{2} (n - 1)^{3}}{2 n - 1} \right).
\end{equation}
When $\Delta=n+1$, then we confirm an upper bound as
\begin{equation}~\label{eqqsigmauppern1}
   n (n + 1) \left( 1 + \frac{2 (n - 1) n^{2} (n + 1)^{2}}{2 n + 1} \right). 
\end{equation}
Actually, both of bounds~(\ref{eqqsigmalowern2}),(\ref{eqqsigmauppern1}) holds to $\sigma(P_n)-M_1(P_n)>0$. Then, 
\begin{equation}~\label{sigmaboundnm1}
     \sigma(P_n) \geq M_1(P_n)  - \frac{2mn^2(\Delta - 1)+\delta-1}{n+\Delta}.
\end{equation}
Finally, by both discussing the two cases~(\ref{sigmaboundn1}),(\ref{sigmaboundnm1}) when the tree is a star tree or a path and in both cases the tree is regular, we have shown all these bounds for the Sigma index and the first Zagreb index, where $\sigma(T)-M_1(T)>0$. As desire.
\end{proof}

\end{document}
